% =========================================================
% epsilon characterization of the supremum
% =========================================================

\subsection{$\varepsilon$-Characterization of Supremum}




% =========================================================
% Epsilon Characterizations of Supremum and Infimum
% House style
% =========================================================

\begin{tcolorbox}[colback=propbox, colframe=propborder, arc=2pt,
  left=6pt, right=6pt, top=4pt, bottom=4pt,
  title={\small\textbf{Definition (Epsilon Characterization of the Supremum)}},
  fonttitle=\small\bfseries]
\begin{definition}
\label{def:epsilon-characterization-of-supremum}
Let $A \subseteq \mathbb{R}$ and let $s \in \mathbb{R}$. The element $s$ is said to satisfy
the \emph{epsilon characterization of the supremum of $A$} if $s$ is an upper bound of $A$
and for every $\varepsilon > 0$ there exists $a \in A$ such that
\[
s-\varepsilon < a.
\]
\end{definition}
\end{tcolorbox}

\begin{remark*}[Standard quantified statement]
\[
\begin{aligned}
&s \text{ satisfies the epsilon characterization of } \sup A \\
&\quad\Longleftrightarrow\quad
 \bigl(\forall x \in A \;(x \le s)\bigr) \\
&\qquad\qquad\land\;
 \bigl(\forall \varepsilon > 0 \;\exists a \in A \;(s-\varepsilon < a)\bigr).
\end{aligned}
\]
\end{remark*}

\begin{remark*}[Definition predicate reading]
\[
\begin{aligned}
&EpsilonCharacterizationOfSupremum(s,A) \\
&\quad\Longleftrightarrow\quad
 \underbrace{UpperBound(s,A)}_{\text{$s$ is an upper bound of $A$}} \\
&\qquad\qquad\land\;
 \underbrace{\forall \varepsilon > 0 \;\exists a \in A \;(s-\varepsilon < a)}_{\text{$A$ comes arbitrarily close to $s$ from below}}.
\end{aligned}
\]
\end{remark*}

\begin{remark*}[Negated quantified statement]
\[
\begin{aligned}
&s \text{ does not satisfy the epsilon characterization of } \sup A \\
&\quad\Longleftrightarrow\quad
 \bigl(\exists x \in A \;(s < x)\bigr) \\
&\qquad\qquad\lor\;
 \bigl(\exists \varepsilon_0 > 0 \;\forall a \in A \;(a \le s-\varepsilon_0)\bigr).
\end{aligned}
\]
\end{remark*}

\begin{remark*}[Failure mode decomposition]
\[
\begin{aligned}
&s \text{ does not satisfy the epsilon characterization of } \sup A \\
&\quad\Longleftrightarrow\quad
 \underbrace{\exists x \in A \;(s < x)}_{\text{$s$ fails to be an upper bound}} \\
&\qquad\qquad\lor\;
 \underbrace{\exists \varepsilon_0 > 0 \;\forall a \in A \;(a \le s-\varepsilon_0)}_{\text{$A$ stays a positive distance below $s$}}.
\end{aligned}
\]
\end{remark*}

\begin{remark*}[Failure modes]\hfill
\begin{itemize}
  \item \textbf{Failure to be an upper bound.} There exists an element of $A$ strictly larger than $s$.
  \item \textbf{Failure of approximation from below.} There exists a positive gap below $s$ that contains no points of $A$ close enough to reach $s$.
\end{itemize}
\end{remark*}

\begin{remark*}[Interpretation]
This characterization says that $s$ sits above all elements of $A$, but not with any positive
slack. No matter how small a positive distance $\varepsilon$ is chosen, some point of $A$
lies within $\varepsilon$ of $s$ from below.

The important point is that the second condition rules out the possibility that $s$ is merely
an upper bound that sits strictly above the set. It forces $A$ to press upward arbitrarily
close to $s$.
\end{remark*}



















\input{volume-iii/analysis/bounding/notes/bounds-extremals/figure-epsilon-characterization-supremum}

\begin{tcolorbox}[colback=propbox, colframe=propborder, arc=2pt,
  left=6pt, right=6pt, top=4pt, bottom=4pt,
  title={\small\textbf{Definition (Epsilon Characterization of the Infimum)}},
  fonttitle=\small\bfseries]
\begin{definition}
\label{def:epsilon-characterization-of-infimum}
Let $A \subseteq \mathbb{R}$ and let $t \in \mathbb{R}$. The element $t$ is said to satisfy
the \emph{epsilon characterization of the infimum of $A$} if $t$ is a lower bound of $A$
and for every $\varepsilon > 0$ there exists $a \in A$ such that
\[
a < t+\varepsilon.
\]
\end{definition}
\end{tcolorbox}

\begin{remark*}[Standard quantified statement]
\[
\begin{aligned}
&t \text{ satisfies the epsilon characterization of } \inf A \\
&\quad\Longleftrightarrow\quad
 \bigl(\forall x \in A \;(t \le x)\bigr) \\
&\qquad\qquad\land\;
 \bigl(\forall \varepsilon > 0 \;\exists a \in A \;(a < t+\varepsilon)\bigr).
\end{aligned}
\]
\end{remark*}

\begin{remark*}[Definition predicate reading]
\[
\begin{aligned}
&EpsilonCharacterizationOfInfimum(t,A) \\
&\quad\Longleftrightarrow\quad
 \underbrace{LowerBound(t,A)}_{\text{$t$ is a lower bound of $A$}} \\
&\qquad\qquad\land\;
 \underbrace{\forall \varepsilon > 0 \;\exists a \in A \;(a < t+\varepsilon)}_{\text{$A$ comes arbitrarily close to $t$ from above}}.
\end{aligned}
\]
\end{remark*}

\begin{remark*}[Negated quantified statement]
\[
\begin{aligned}
&t \text{ does not satisfy the epsilon characterization of } \inf A \\
&\quad\Longleftrightarrow\quad
 \bigl(\exists x \in A \;(x < t)\bigr) \\
&\qquad\qquad\lor\;
 \bigl(\exists \varepsilon_0 > 0 \;\forall a \in A \;(t+\varepsilon_0 \le a)\bigr).
\end{aligned}
\]
\end{remark*}

\begin{remark*}[Failure mode decomposition]
\[
\begin{aligned}
&t \text{ does not satisfy the epsilon characterization of } \inf A \\
&\quad\Longleftrightarrow\quad
 \underbrace{\exists x \in A \;(x < t)}_{\text{$t$ fails to be a lower bound}} \\
&\qquad\qquad\lor\;
 \underbrace{\exists \varepsilon_0 > 0 \;\forall a \in A \;(t+\varepsilon_0 \le a)}_{\text{$A$ stays a positive distance above $t$}}.
\end{aligned}
\]
\end{remark*}

\begin{remark*}[Failure modes]\hfill
\begin{itemize}
  \item \textbf{Failure to be a lower bound.} There exists an element of $A$ strictly smaller than $t$.
  \item \textbf{Failure of approximation from above.} There exists a positive gap above $t$ that contains no points of $A$ close enough to reach $t$.
\end{itemize}
\end{remark*}

\begin{remark*}[Interpretation]
This characterization says that $t$ sits below all elements of $A$, but not with any positive
slack. No matter how small a positive distance $\varepsilon$ is chosen, some point of $A$
lies within $\varepsilon$ of $t$ from above.

The important point is that the second condition rules out the possibility that $t$ is merely
a lower bound that sits strictly below the set. It forces $A$ to press downward arbitrarily
close to $t$.
\end{remark*}



